%=========================================================
% C++ Chapter 11 Lab Handout (Verbatim Korean-safe)
% Topic: C++ I/O System (Streams, Formatting, Manipulators, << >> Overloads)
% Compile with XeLaTeX
%=========================================================
\documentclass[11pt,a4paper]{article}

\usepackage[margin=1in]{geometry}
\usepackage{fontspec}
\usepackage{kotex}
\usepackage{hyperref}
\usepackage{enumitem}
\usepackage{fancyhdr}
\usepackage{titlesec}
\usepackage{xcolor}
\usepackage{tcolorbox}
\usepackage{fvextra}

% ---------- Fonts ----------
% If D2Coding isn't available, try: NanumGothicCoding, Noto Sans Mono CJK KR
\setmonofont{D2Coding}[Scale=MatchLowercase]

% ---------- Colors ----------
\definecolor{CodeBg}{RGB}{248,248,248}
\definecolor{CodeFrame}{RGB}{220,220,220}
\definecolor{Accent}{RGB}{40,80,160}

% ---------- Header ----------
\pagestyle{fancy}
\fancyhf{}
\lhead{\textbf{C++ 11장 실습}}
\rhead{\textbf{C++ 입출력 시스템}}
\cfoot{\thepage}

\titleformat{\section}{\Large\bfseries\color{Accent}}{\thesection}{0.6em}{}
\titleformat{\subsection}{\large\bfseries}{\thesubsection}{0.6em}{}

% ---------- Boxes ----------
\newtcolorbox{GoalBox}{
  colback=CodeBg,
  colframe=CodeFrame,
  arc=3pt,
  boxrule=0.8pt,
  left=10pt,right=10pt,top=8pt,bottom=8pt
}
\newcommand{\filetag}[1]{\texttt{\color{Accent}#1}}

% ---------- Code block environment (Unicode-safe) ----------
% IMPORTANT: do NOT set commandchars, so C++ braces {} print correctly.
\DefineVerbatimEnvironment{cppcode}{Verbatim}{
  fontsize=\small,
  frame=single,
  framerule=0.6pt,
  rulecolor=\color{CodeFrame},
  numbers=left,
  numbersep=8pt,
  baselinestretch=1,
  breaklines=true,
  breaksymbolleft=,
  breaksymbolright=,
  xleftmargin=1.5em
}

\begin{document}

\begin{center}
  {\LARGE \textbf{C++ 11장 실습: 입출력 시스템}}\\
  \vspace{6pt}
  {\large (stream/buffer, istream/ostream members, formatted I/O, manipulators, \texttt{<< >>} overloading)}
\end{center}

\vspace{8pt}

\section*{학습 목표}
\begin{GoalBox}
\begin{enumerate}[leftmargin=18pt]
  \item 스트림(stream) 입출력과 버퍼의 의미를 이해한다. % p.3~8
  \item ostream의 \texttt{put()}, \texttt{write()}, \texttt{flush()}를 사용할 수 있다. % p.15~16
  \item istream의 \texttt{get()}, \texttt{getline()}, \texttt{gcount()}, \texttt{ignore()}를 사용할 수 있다. % p.17~24
  \item 포맷 입출력(플래그/함수/조작자)을 이해하고 사용할 수 있다. % p.25~36
  \item 삽입/추출 연산자(\texttt{<<}, \texttt{>>})의 실행 과정을 이해한다. % p.37~46
  \item 사용자 정의 \texttt{<<}, \texttt{>>}, 조작자를 작성할 수 있다. % p.39~50
\end{enumerate}
\end{GoalBox}

\section*{실습 운영 규칙}
\begin{itemize}[leftmargin=18pt]
  \item 모든 예제는 \texttt{main()} 포함, \textbf{독립 실행}입니다.
  \item 입력 예제는 \textbf{버퍼에 남는 '\textbackslash n'} 문제를 일부러 다룹니다. (핵심!)
  \item 포맷 예제는 \textbf{cout 상태(state)가 유지}된다는 점을 주의해서 관찰하세요.
\end{itemize}

\tableofcontents
\newpage

%=========================================================
\section{ostream 멤버 함수: put(), write(), flush() (예제 11-1)}
\begin{GoalBox}
\textbf{핵심}
\begin{itemize}[leftmargin=18pt]
  \item \texttt{put(char)}: 문자 1개 출력. 정수(ASCII 코드)를 넣어도 됨(예: 33='!'). % p.16
  \item \texttt{write(char*, n)}: 배열에서 n개 문자 출력(널문자까지 자동 출력 아님). % p.15~16
  \item \texttt{flush()}: 출력 버퍼를 즉시 비움(강제 출력). % p.7, p.15
\end{itemize}
\end{GoalBox}

\noindent\filetag{L1\_ostream\_put\_write.cpp}
\begin{cppcode}
#include <iostream>
using namespace std;

int main() {
    // "Hi!" 출력 후 줄바꿈
    cout.put('H');
    cout.put('i');
    cout.put(33);     // ASCII 33 == '!'
    cout.put('\n');

    // put()은 체이닝 가능(ostream& 리턴)
    cout.put('C').put('+').put('+').put(' ');

    char str[] = "I love programming";
    cout.write(str, 6); // "I love" 6글자만 출력
    cout.put('\n');

    // flush 강제 출력(보통은 '\n' 또는 버퍼가 찰 때 출력됨)
    cout << "[flush demo]" << flush;
}
\end{cppcode}

%=========================================================
\section{istream get(): 한 글자씩 읽기(EOF/Enter 처리) (예제 11-2)}
\begin{GoalBox}
\textbf{핵심}
\begin{itemize}[leftmargin=18pt]
  \item \texttt{int get()}는 EOF(-1)과 비교해야 해서 int로 받는 패턴이 흔함. % p.17
  \item \texttt{istream\& get(char\&)}는 문자 1개를 참조로 채움.
  \item '\textbackslash n'(Enter)을 만나면 한 줄 읽기 종료로 처리 가능.
\end{itemize}
\end{GoalBox}

\noindent\filetag{L2\_istream\_get\_line.cpp}
\begin{cppcode}
#include <iostream>
using namespace std;

void get1() {
    cout << "cin.get()로 <Enter>까지 입력 받고 출력합니다>> ";
    int ch; // EOF 비교를 위해 int
    while ((ch = cin.get()) != EOF) {
        cout.put(ch);
        if (ch == '\n') break; // Enter에서 중단
    }
}

void get2() {
    cout << "cin.get(char&)로 <Enter>까지 입력 받고 출력합니다>> ";
    char ch;
    while (true) {
        cin.get(ch);
        if (cin.eof()) break;
        cout.put(ch);
        if (ch == '\n') break;
    }
}

int main() {
    get1();
    get2();
}
\end{cppcode}

%=========================================================
\section{get(char*, n)과 '\\textbackslash n' 버퍼 함정 + ignore() (예제 11-3)}
\begin{GoalBox}
\textbf{핵심}
\begin{itemize}[leftmargin=18pt]
  \item \texttt{cin.get(buf,n)}는 '\textbackslash n'을 만나면 \textbf{읽기를 멈추고}, '\textbackslash n'은 \textbf{버퍼에 남겨둠}. % p.20~21
  \item 그래서 다음 입력이 바로 '\textbackslash n'을 만나 즉시 종료될 수 있음(무한 루프 원인).
  \item 해결: \texttt{cin.ignore(1);} 또는 \texttt{cin.get();}로 '\textbackslash n' 제거.
\end{itemize}
\end{GoalBox}

\noindent\filetag{L3\_get\_charptr\_ignore.cpp}
\begin{cppcode}
#include <iostream>
#include <cstring>
using namespace std;

int main() {
    char cmd[80];

    cout << "cin.get(char*, int)로 문자열을 읽습니다.\n";
    while (true) {
        cout << "종료하려면 exit을 입력하세요 >> ";
        cin.get(cmd, 80); // '\n' 전까지 읽고, '\n'은 버퍼에 남김

        if (strcmp(cmd, "exit") == 0) {
            cout << "프로그램을 종료합니다....\n";
            return 0;
        }

        // (핵심) 버퍼에 남은 '\n' 제거. 이 줄 없으면 다음 루프에서 바로 리턴될 수 있음.
        cin.ignore(1);
    }
}
\end{cppcode}

%=========================================================
\section{getline(): 공백 포함 한 줄 읽기 (예제 11-4)}
\begin{GoalBox}
\textbf{핵심}
\begin{itemize}[leftmargin=18pt]
  \item \texttt{cin.getline(line,n)}는 '\textbackslash n'을 만나면 종료하고, '\textbackslash n'은 \textbf{버퍼에서 제거}되며 배열에는 넣지 않음. % p.22~23
  \item 공백 포함 문자열을 안정적으로 다룰 때 자주 사용.
\end{itemize}
\end{GoalBox}

\noindent\filetag{L4\_getline\_echo.cpp}
\begin{cppcode}
#include <iostream>
#include <cstring>
using namespace std;

int main() {
    char line[80];
    cout << "cin.getline() 함수로 라인을 읽습니다.\n";
    cout << "exit를 입력하면 루프가 끝납니다.\n";

    int no = 1;
    while (true) {
        cout << "라인 " << no << " >> ";
        cin.getline(line, 80);
        if (strcmp(line, "exit") == 0) break;
        cout << "echo --> " << line << "\n";
        no++;
    }
}
\end{cppcode}

%=========================================================
\section{gcount()와 ignore(): 읽은 문자 수/건너뛰기 (슬라이드 p.24)}
\begin{GoalBox}
\textbf{핵심}
\begin{itemize}[leftmargin=18pt]
  \item \texttt{cin.gcount()}는 \textbf{직전에 실행한 입력 함수가 읽은 문자 수}를 리턴. % p.24
  \item \texttt{cin.ignore(n)}: n개 문자 제거
  \item \texttt{cin.ignore(n, delim)}: delim을 만나면 그 문자까지 제거하고 종료
\end{itemize}
\end{GoalBox}

\noindent\filetag{L5\_gcount\_ignore.cpp}
\begin{cppcode}
#include <iostream>
using namespace std;

int main() {
    char line[80];
    cout << "문장을 입력하세요>> ";
    cin.getline(line, 80);

    int n = (int)cin.gcount(); // '\0'은 포함되지 않는다고 가정(환경에 따라 다를 수 있음)
    cout << "getline이 읽은 문자 수(gcount) = " << n << "\n";
    cout << "입력 내용 = [" << line << "]\n";

    cout << "이제 입력 스트림에서 5문자 건너뜁니다. (ignore)\n";
    cin.ignore(5);

    cout << "다음 입력(한 단어)>> ";
    string w;
    cin >> w;
    cout << "w = " << w << "\n";
}
\end{cppcode}

%=========================================================
\section{포맷 플래그: setf()/unsetf() (예제 11-5)}
\begin{GoalBox}
\textbf{핵심}
\begin{itemize}[leftmargin=18pt]
  \item \texttt{cout.setf(ios::hex)} 등으로 출력 형식을 설정. % p.28~29
  \item \texttt{unsetf(ios::dec)}처럼 플래그를 해제할 수 있음.
  \item 플래그는 \textbf{cout 상태로 유지}되므로 다음 출력에도 영향을 줌.
\end{itemize}
\end{GoalBox}

\noindent\filetag{L6\_setf\_unsetf.cpp}
\begin{cppcode}
#include <iostream>
using namespace std;

int main() {
    cout << 30 << "\n";            // 기본: 10진수
    cout.unsetf(ios::dec);         // 10진수 해제
    cout.setf(ios::hex);           // 16진수
    cout << 30 << "\n";            // 1e

    cout.setf(ios::showbase);      // 0x 출력
    cout << 30 << "\n";            // 0x1e

    cout.setf(ios::uppercase);     // A~F 대문자
    cout << 30 << "\n";            // 0X1E

    cout.setf(ios::dec | ios::showpoint);
    cout << 23.5 << "\n";          // 23.5000...

    cout.setf(ios::scientific);
    cout << 23.5 << "\n";          // 2.350000E+001

    cout.setf(ios::showpos);
    cout << 23.5 << "\n";          // +2.350000E+001
}
\end{cppcode}

%=========================================================
\section{포맷 함수: width(), fill(), precision() (예제 11-6)}
\begin{GoalBox}
\textbf{핵심}
\begin{itemize}[leftmargin=18pt]
  \item \texttt{width(n)}는 \textbf{다음 한 번의 출력에만} 적용.
  \item \texttt{fill(ch)}은 이후 출력에 계속 영향을 줄 수 있음(상태 유지).
  \item \texttt{precision(p)}은 유효숫자/소수 표현에 영향(상태 유지). % p.30~31
\end{itemize}
\end{GoalBox}

\noindent\filetag{L7\_width\_fill\_precision.cpp}
\begin{cppcode}
#include <iostream>
using namespace std;

void showWidth() {
    cout.width(10);
    cout << "Hello" << "\n";

    cout.width(5);
    cout << 12 << "\n";

    cout << '%';
    cout.width(10);
    cout << "Korea/" << "Seoul/" << "City" << "\n";
}

int main() {
    showWidth();
    cout << "\n";

    cout.fill('^');   // 상태 유지
    showWidth();
    cout << "\n";

    cout.precision(5);
    cout << 11.0/3.0 << "\n";
}
\end{cppcode}

%=========================================================
\section{조작자: 매개변수 없는 조작자/있는 조작자 (예제 11-7, 11-8)}
\begin{GoalBox}
\textbf{핵심}
\begin{itemize}[leftmargin=18pt]
  \item 조작자는 \textbf{함수}이며, 항상 \texttt{<<} 또는 \texttt{>>}와 함께 사용. % p.32
  \item 매개변수 없는 조작자: \texttt{hex}, \texttt{dec}, \texttt{showbase}, \texttt{boolalpha} 등. % p.33~35
  \item 매개변수 있는 조작자: \texttt{setw}, \texttt{setfill}, \texttt{setprecision} 등(\texttt{<iomanip>} 필요). % p.34~36
\end{itemize}
\end{GoalBox}

\noindent\filetag{L8\_manipulators\_basic.cpp}
\begin{cppcode}
#include <iostream>
#include <iomanip>
using namespace std;

int main() {
    cout << hex << showbase << 30 << "\n";
    cout << dec << showpos << 100 << "\n";

    cout << true << ' ' << false << "\n";         // 1 0 처럼 출력될 수 있음
    cout << boolalpha << true << ' ' << false << "\n"; // true false

    cout << setw(10) << setfill('^') << "Hello" << "\n";

    cout << showbase;
    cout << setw(8) << "Number" << setw(10) << "Octal" << setw(10) << "Hexa" << "\n";
    for (int i = 0; i < 50; i += 5) {
        cout << setw(8) << setfill('.') << dec << i;
        cout << setw(10) << setfill(' ') << oct << i;
        cout << setw(10) << setfill(' ') << hex << i << "\n";
    }
}
\end{cppcode}

%=========================================================
\section{삽입 연산자(<<) 오버로딩: Point 출력 (예제 11-9)}
\begin{GoalBox}
\textbf{핵심}
\begin{itemize}[leftmargin=18pt]
  \item \texttt{cout << p}가 되려면 \texttt{operator<<(ostream\&, Point)}를 제공해야 함. % p.39~41
  \item 체이닝을 위해 \textbf{ostream\&를 리턴}한다. (7장과 연결)
  \item private 멤버 접근이 필요하면 friend 사용.
\end{itemize}
\end{GoalBox}

\noindent\filetag{L9\_Point\_ostream.cpp}
\begin{cppcode}
#include <iostream>
using namespace std;

class Point {
    int x, y;
public:
    Point(int x=0, int y=0) : x(x), y(y) {}
    friend ostream& operator<<(ostream& os, const Point& a);
};

ostream& operator<<(ostream& os, const Point& a) {
    os << "(" << a.x << "," << a.y << ")";
    return os;
}

int main() {
    Point p(3,4);
    cout << p << "\n";
    Point q(1,100), r(2,200);
    cout << q << r << "\n"; // 연속 출력(체이닝)
}
\end{cppcode}

%=========================================================
\section{<< 오버로딩: Book 출력 (예제 11-10)}
\noindent\filetag{L10\_Book\_ostream.cpp}
\begin{cppcode}
#include <iostream>
#include <string>
using namespace std;

class Book {
    string title, press, author;
public:
    Book(string title="", string press="", string author="")
        : title(title), press(press), author(author) {}
    friend ostream& operator<<(ostream& os, const Book& b);
};

ostream& operator<<(ostream& os, const Book& b) {
    os << b.title << "," << b.press << "," << b.author;
    return os;
}

int main() {
    Book book("소유냐 존재냐", "한국출판사", "에리히프롬");
    cout << book << "\n";
}
\end{cppcode}

%=========================================================
\section{추출 연산자(>>) 오버로딩: Point 입력 + 출력 (예제 11-11)}
\begin{GoalBox}
\textbf{핵심}
\begin{itemize}[leftmargin=18pt]
  \item \texttt{cin >> p}가 되려면 \texttt{operator>>(istream\&, Point\&)}를 제공해야 함. % p.44~46
  \item 입력 대상은 변경되므로 \textbf{참조(Point\&)}로 받는다.
  \item 체이닝을 위해 \textbf{istream\&를 리턴}한다.
\end{itemize}
\end{GoalBox}

\noindent\filetag{L11\_Point\_istream\_ostream.cpp}
\begin{cppcode}
#include <iostream>
using namespace std;

class Point {
    int x, y;
public:
    Point(int x=0, int y=0) : x(x), y(y) {}
    friend istream& operator>>(istream& ins, Point& a);
    friend ostream& operator<<(ostream& os, const Point& a);
};

istream& operator>>(istream& ins, Point& a) {
    cout << "x 좌표>>";
    ins >> a.x;
    cout << "y 좌표>>";
    ins >> a.y;
    return ins;
}

ostream& operator<<(ostream& os, const Point& a) {
    os << "(" << a.x << "," << a.y << ")";
    return os;
}

int main() {
    Point p;
    cin >> p;
    cout << p << "\n";
}
\end{cppcode}

%=========================================================
\section{사용자 정의 조작자: beep, rightarrow, fivestar (예제 11-12)}
\begin{GoalBox}
\textbf{핵심}
\begin{itemize}[leftmargin=18pt]
  \item 매개변수 없는 조작자 형태: \texttt{ostream\&(ostream\&)}. % p.49
  \item \texttt{cout << beep << endl;} 처럼 사용 가능.
\end{itemize}
\end{GoalBox}

\noindent\filetag{L12\_custom\_manipulators.cpp}
\begin{cppcode}
#include <iostream>
using namespace std;

ostream& fivestar(ostream& outs) { return outs << "*****"; }
ostream& rightarrow(ostream& outs) { return outs << "---->"; }
ostream& beep(ostream& outs) { return outs << '\a'; } // 시스템 beep

int main() {
    cout << "벨이 울립니다" << beep << endl;
    cout << "C" << rightarrow << "C++" << rightarrow << "Java" << endl;
    cout << "Visual" << fivestar << "C++" << endl;
}
\end{cppcode}

%=========================================================
\section{istream 조작자: 질문 출력 후 입력 받기 (예제 11-13)}
\begin{GoalBox}
\textbf{핵심}
\begin{itemize}[leftmargin=18pt]
  \item 입력 조작자 형태: \texttt{istream\&(istream\&)}. % p.50
  \item \texttt{cin >> question >> answer;} 처럼 ``입력 전 훅(hook)''으로 활용 가능.
\end{itemize}
\end{GoalBox}

\noindent\filetag{L13\_istream\_manipulator\_question.cpp}
\begin{cppcode}
#include <iostream>
#include <string>
using namespace std;

istream& question(istream& ins) {
    cout << "거울아 거울아 누가 제일 예쁘니?";
    return ins;
}

int main() {
    string answer;
    cin >> question >> answer;
    cout << "세상에서 제일 예쁜 사람은 " << answer << "입니다." << endl;
}
\end{cppcode}

\section*{마무리 체크리스트}
\begin{GoalBox}
\begin{itemize}[leftmargin=18pt]
  \item 스트림 버퍼가 왜 필요한지(입력 수정/출력 효율)를 설명할 수 있는가? % p.5~7
  \item \texttt{get()/getline()}에서 '\textbackslash n'이 버퍼에 남는 경우와 제거되는 경우를 구분할 수 있는가? % p.20~23
  \item \texttt{ignore()}를 언제 써야 하는지 설명할 수 있는가? % p.20~24
  \item setf/unsetf, width/fill/precision의 ``상태 유지'' 차이를 설명할 수 있는가? % p.28~31
  \item 매개변수 없는/있는 조작자를 구분할 수 있는가? % p.32~36
  \item \texttt{<<} / \texttt{>>} 오버로딩에서 왜 \texttt{ostream\&}/\texttt{istream\&}를 리턴해야 하는가? % p.37~46
  \item 사용자 정의 조작자의 함수 시그니처를 쓸 수 있는가? % p.49~50
\end{itemize}
\end{GoalBox}

\end{document}
